\begin{statement}[C.8.2; Proved: an explicit local maximum on a larger disc]
If $|\tau-i|\le1/50$ and $\wp_\tau(z)=0$, then
$N(z,\tau)\le\sqrt2$, with equality exactly when $\tau=i$.
\end{statement}
\begin{proof}
Write $\delta=\tau-i$, $h=z-(1-\tau)/2$ and
\[
 G(h,\tau)=\wp_\tau((1-\tau)/2+h),\qquad
 H(h,\tau)=\log\frac{\vartheta_3(h,\tau)}{\eta(\tau)}.
\]
We choose the holomorphic logarithm whose value at $(0,i)$ is
$\log\sqrt2$. The half-period theta identity gives
\begin{equation}\label{eq:ca-H}
 \log N(z,\tau)=\Re H(h,\tau)-\frac{\pi(\Im h)^2}{T}.
\end{equation}
Both $G$ and $H$ are even in $h$. On the polydisc
$|h|\le1/4$, $|\delta|\le1/10$, they satisfy
\begin{equation}\label{eq:ca-GHbounds}
 |G|<30,\qquad |H|<2/3.
\end{equation}
Here are sufficient series estimates, including their margins.
Put
\[
 A_k=\sum_{n\ge1}n^k e^{-9\pi n^2/10}\cosh(\pi n/2),
 \quad a_0=2A_0,\quad a_1=4\pi A_1,\quad a_2=8\pi^2A_2.
\]
Then $a_0<3/10$, $a_1<19/10$, $a_2<12$.
The first two summands suffice to verify these inequalities: for
$n\ge3$, each of the two exponentials in the hyperbolic cosine
has successive-term ratio at most $e^{-29\pi/5}$, and the
factor $n^k$ increases this ratio by at most $(4/3)^2$.
We use \eqref{eq:ca-exp} for the first terms and this geometric
bound for the tail. With $r=e^{-9\pi/5}$, the identities
\[
 G=-\partial_h^2\log\vartheta_3-\frac{\pi^2}{3}E_2,
 \qquad
 |E_2|\le1+\frac{24r}{(1-r)^3}
\]
give
\[
 |G|\le\frac{a_2}{1-a_0}+\frac{a_1^2}{(1-a_0)^2}
       +\frac{\pi^2}{3}\left(1+\frac{24r}{(1-r)^3}\right)<30.
\]
Also
\[
 |H|\le-\log(1-a_0)+\frac{11\pi}{120}
                    +\frac{r}{(1-r)^2}<\frac23.
\]
The logarithm is well defined since $|\vartheta_3-1|<3/10$;
the displayed eta product fixes its branch. A geometric bound
for $-\log(1-x)=\sum_{n\ge1}x^n/n$ verifies the last estimate.

Let
\[
 e(\tau)=\frac{\pi^2}{3}(\vartheta_2(0,\tau)^4-
                              \vartheta_4(0,\tau)^4),\qquad
 E(\tau)=\frac{\pi^2}{3}E_2(\tau),\qquad
 g=\frac{g_2(i)}4.
\]
On $|\delta|\le1/50$ the following loose bounds hold:
\begin{equation}\label{eq:ca-derivatives}
 |e'|<14,\quad |e''|<60,\quad |g_2'|<500,
 \quad |E'|<2,\quad |E_6''|<100,
 \qquad 47<g<95/2.
\end{equation}
For completeness these are scalar Fourier-series bounds, not
interval computations. At $T\ge49/50$, upper bounds for
$|\vartheta_j|,|\vartheta_j'|,|\vartheta_j''|$ are
\[
\begin{array}{c|ccc}
j&0\text{ derivatives}&1\text{ derivative}&2\text{ derivatives}\\ \hline
2&929/1000&742/1000&670/1000\\
4&1093/1000&290/1000&911/1000.
\end{array}
\]
They follow by differentiating the absolutely convergent theta
series, taking absolute values, retaining two terms, and using
the same successive-term bound. Differentiating the fourth
powers proves the first two bounds in \eqref{eq:ca-derivatives}.
For $r\le1/470$, the remaining derivative majorants are
\begin{align*}
 |g_2'|&\le640\pi^5\frac{r(1+11r+11r^2+r^3)}{(1-r)^7}<500,\\
 |E'|&\le16\pi^3\frac{r(1+r)}{(1-r)^5}<2,\\
 |E_6''|&\le2016\pi^2
 \frac{r(1+r)(1+120r+1191r^2+2416r^3+1191r^4+120r^5+r^6)}
 {(1-r)^{11}}<100.
\end{align*}
These follow by summing $\sum n^k r^n$ and bounding
$(1-r^n)^{-1}$ by $(1-r)^{-1}$. Substitution of $r=1/470$
and $\pi<355/113$ verifies the rational inequalities.
Finally $g=\pi^4E_4(i)/3$ and
$E_4=1+240\sum n^3q^n/(1-q^n)$ give $47<g<95/2$.
The identity of K3, checked by differentiating $e$ and using
the theta heat equation, reads
\[
 e(i)=0,\qquad e'(i)=id,\qquad d=\frac{2g}{3\pi}>99/10.
\]

The differential equation $\wp''=6\wp^2-g_2/2$ determines the
first coefficients in the even expansion:
\[
 G(h,\tau)=e+a h^2+ea h^4+\mathcal R_G,
 \qquad a=3e^2-g_2/4.
\]
Consequently, on $|\delta|\le1/50$,
\[
 |e|\le7/25,\quad |a+g|\le\frac52+3(7/25)^2=\frac{3419}{1250},
 \quad |a|\le\frac{62794}{1250}.
\]
Cauchy's estimate from \eqref{eq:ca-GHbounds} yields
\[
 |\mathcal R_G|\le
 \frac{30(4|h|)^6}{1-16|h|^2}.
\]
On $|h|=1/10$ the term $ah^2$ has modulus at least
$55331/125000>0.4426$, whereas the sum of the other terms is at most
\[
 \frac7{25}+\frac7{25}\frac{62794}{1250}\frac1{10000}
       +\frac{30(4/25)^3}{1-4/25}<0.4277.
\]
Rouch\'e's theorem puts two zeros, with multiplicity, in
$|h|<1/10$. The disc injects into the torus and $\wp$ has exactly
two zeros on the torus, so it contains both of them.

Write $G=e+h^2Q$. The preceding estimates give, for $|h|\le1/10$,
\[
 |Q+g|\le
 \frac{3419}{1250}+\frac7{25}\frac{62794}{1250}\frac1{100}
       +\frac{30(4/25)^3}{(1-4/25)/100}<\frac{88}{5}.
\]
Taylor's formula and $|e''|<60$ give
$|e|\ge(99/10-3/5)|\delta|=(93/10)|\delta|$.
At either zero, therefore,
\begin{equation}\label{eq:ca-delta-h}
 |\delta|<7|h|^2,
\end{equation}
because $(95/2+88/5)/(93/10)=7$.

Expand
$H=H_0(\tau)+H_2(\tau)h^2+H_4(\tau)h^4+\mathcal R_H$.
The heat equation and the identity for $G$ give
\[
 H_0'=\frac{ie}{4\pi},\quad
 H_2=-\frac{e+E}{2},\quad H_4=-\frac a{12}.
\]
Thus $H_0(i)=\log\sqrt2$, $H_0'(i)=0$,
$|H_0''|<6/5$, $H_2(i)=-\pi/2$, $|H_2'|<8$, and
$|H_4|<21/5$. Cauchy's estimate gives
\[
 |\mathcal R_H|\le
 \frac{(2/3)(4|h|)^6}{1-16|h|^2}
 \le\frac{2048}{63}|h|^4
 \quad (|h|\le1/10).
\]
Combining \eqref{eq:ca-H} and \eqref{eq:ca-delta-h}, we obtain
\begin{align*}
 \log N-\log\sqrt2
 &\le-\frac\pi2\Re(h^2)-\pi(\Im h)^2+C|h|^4\\
 &=-\frac\pi2|h|^2+C|h|^4,\\
 C&=\frac{147}{5}+56+\frac{21}{5}
                 +\frac{2048}{63}+\frac{1100}{49}<150.
\end{align*}
The final term of $C$ bounds the change of $1/T$ from $1$:
$7\pi/(1-1/50)<1100/49$. Since $150|h|^2<3/2<\pi/2$,
the result is strict unless $h=0$. Equation
\eqref{eq:ca-delta-h} then gives $\delta=0$, where equality
holds by the checked isogeny computation at $i$.
\end{proof}
